\noindent\textbf{Extended summary}\par\smallskip
High-performance centres invest in force plates and test athletes routinely, yet decisions about those athletes usually rest on a small number of vendor-computed values, such as jump height and peak force. Athletes with similar values are then treated as though they had jumped in the same way. This study asks whether that inference holds and how much of the recorded information those values leave unused.

The analysis used \nTrials{} countermovement jumps from \nAth{} male football players tested during routine monitoring at a private sports physiotherapy and performance centre in India. Each force--time curve was normalised to body weight; the take-off and landing phases were time-normalised separately and concatenated, and principal component analysis identified the main modes of variation. The dispersion of curve shapes among jumps matched on jump height, or on relative peak force, was compared with that of the whole sample, and the share of each mode explained by four vendor metrics, age and body mass was estimated.

Four components explained \cumFour\% of curve variance, and all four lay almost entirely in the landing phase; the first described how quickly the athlete arrested the fall. Jumps matched on jump height were scarcely more alike in shape than jumps in general, and matching on relative peak force had no discernible effect. The conventional metrics, age and body mass explained \rtwoRange\% of each component (\cvRange\% out of sample). Between-athlete differences accounted for \betweenRange\% of component variance, indicating that much of the variation was trial-specific. With a fixed 1-s landing window, a take-off component emerged that the metrics largely explained.

For those who manage sports-science programmes, the landing is recorded in every jump at no additional cost. Simple landing descriptors could be added to routine reports, provided that decisions rest on several trials rather than on a single landing.
